\documentclass[11pt]{article}
\usepackage[T1]{fontenc}
\usepackage{lmodern}
\usepackage{amsmath,amssymb,amsthm}
\usepackage[margin=1in]{geometry}
\usepackage{microtype}
\usepackage[dvipsnames]{xcolor}
\usepackage{hyperref}
\hypersetup{colorlinks=true,linkcolor=MidnightBlue,citecolor=MidnightBlue,urlcolor=MidnightBlue,
 pdftitle={Uniform Neumann control for strictly outward non-gradient flows},pdfauthor={Kaiwen Jin}}
\newtheorem{theorem}{Theorem}[section]
\newtheorem{proposition}[theorem]{Proposition}
\newtheorem{lemma}[theorem]{Lemma}
\theoremstyle{remark}
\newtheorem{remark}[theorem]{Remark}
\newcommand{\eps}{\varepsilon}
\newcommand{\R}{\mathbb R}
\newcommand{\divg}{\operatorname{div}}
\newcommand{\supp}{\operatorname{supp}}
\newcommand{\osc}{\operatorname{osc}}
\newcommand{\dd}{\,dx\,dt}
\newcommand{\ds}{\,dS\,dt}
\allowdisplaybreaks[2]

\title{Uniform Neumann control for strictly outward non-gradient flows}
\author{Kaiwen Jin\\\small\href{mailto:yui.ui.siki@gmail.com}{yui.ui.siki@gmail.com}}
\date{October 12, 2026}

\begin{document}
\maketitle
\begin{abstract}
We prove uniform null controllability of an advection--diffusion equation with
Neumann boundary conditions as the diffusion coefficient tends to zero. The
drift is autonomous, belongs to $W^{2,\infty}$, points strictly outward on the
boundary, and its backward flow carries a neighborhood of every point into
the interior control region within a bounded time. For every sufficiently
large control time, an $L^2$ control drives any $L^2$ initial state to zero
with a norm bound independent of the diffusion coefficient. No gradient
representation of the drift is required. A scalar function prescribed only
through its boundary normal derivative leaves a tangential drift and a
positive Robin coefficient in the adjoint equation. Two Carleman weights
then cancel the uncontrolled tangential boundary terms. Their fixed-time
observation estimate combines with the known flushing dissipation estimate
to give the uniform control bound.
\end{abstract}

\section{Introduction}

Consider an interior control for an advection--diffusion equation on a bounded
domain. Null controllability asks whether this control can drive every initial
state to zero at a prescribed time. When the diffusion coefficient $\eps$ is
small, the equation is close to transport, so the trajectories of the drift
become relevant to the cost of control. The boundary condition is equally
important: Neumann and mixed conditions can produce a cost growing
exponentially in $1/\eps$, even at large times
\cite[Theorems 2.8 and 2.9]{Barcena}. Thus controllability at each fixed $\eps>0$ does
not by itself provide a uniform bound as $\eps\to0$.

For Neumann equations, Et-Tahri, B\'arcena-Petisco, Boutaayamou and Maniar
developed a Carleman calculation for tangential drifts
\cite[Section 6]{Tangential}. Their uniform controllability statement
\cite[Theorem 1.4]{Tangential} assumes backward flushing into an interior
control region. A Lipschitz tangential flow preserves the boundary, however,
so this hypothesis cannot hold for boundary trajectories
\cite[Remark 1.6]{Outward}. Their subsequent work proves uniform large-time
control for strictly outward gradient drifts \cite[Theorem 1.5]{Outward}.
A gradient representation allows
one to remove the first-order term by a scalar conjugation. For a general
outward drift, the same direct Carleman calculation leaves boundary terms
whose tangential part is not controlled by an interior observation
\cite[Remark 1.6]{Outward}. The dissipation estimate under the trajectory
condition, however, already holds for general drifts
\cite[Proposition 4.7(1)]{Outward}.

Here we retain that dissipation estimate and prove the missing fixed-time
observation bound for an autonomous drift. We use a scalar conjugation to
change its normal component on the boundary, while retaining its tangential
component in the equation. This is sufficient for the boundary cancellation;
the drift need not be a gradient.

Let $\Omega\subset\R^d$, $d\ge2$, be a bounded connected smooth domain, let
$\Gamma=\partial\Omega$, and let $\nu$ be the outward unit normal. Let
$\varnothing\ne\omega\Subset\Omega$ be open and
$B\in W^{2,\infty}(\R^d;\R^d)$. We use its $C^{1,1}$ representative and assume
\begin{equation}\label{outward}
 b:=B\cdot\nu\ge c_0>0\quad\hbox{on }\Gamma.
\end{equation}
The flow $F_s$ of $-B$ is globally defined. We assume that there are
$T_0,r_0>0$ such that
\begin{equation}\label{flushing}
 \forall x\in\overline\Omega\quad\exists s_x\in(0,T_0):
 \quad F_{s_x}\bigl(\mathbb B(x,r_0)\bigr)\subset\omega.
\end{equation}
Here $\mathbb B(x,r_0)$ is the open Euclidean ball. This condition means that a
whole neighborhood of each point, rather than only one trajectory, enters
the control region. It is the ball version of the flushing condition used in
\cite{Outward}, with the autonomous backward flow written explicitly.

We study
\begin{equation}\label{state}
 \begin{cases}
 y_t-\eps\Delta y+B\cdot\nabla y=1_\omega h
                  &\text{in }\Omega\times(0,T),\\
 \partial_\nu y=0  &\text{on }\Gamma\times(0,T),\\
 y(0)=y_0         &\text{in }\Omega.
 \end{cases}
\end{equation}
The solution is understood in the variational sense specified in Section
\ref{sec:adjoint}.

\begin{theorem}\label{main}
Under \eqref{outward}--\eqref{flushing}, there exist finite constants
$T_*>0$ and $C>0$, depending on the fixed domain, control region and drift
data, such that for every $T\ge T_*$, every $0<\eps<1$, and every
$y_0\in L^2(\Omega)$, there is an $h\in L^2((0,T)\times\omega)$ for which the
solution of \eqref{state} satisfies
\[
 y(T)=0\quad\text{in }L^2(\Omega),\qquad
 \|h\|_{L^2((0,T)\times\omega)}\le C\|y_0\|_{L^2(\Omega)}.
\]
Neither $T_*$ nor $C$ depends on $\eps$, and $C$ can be chosen independently
of $T\ge T_*$. No gradient condition on $B$ is imposed.
\end{theorem}

In particular, the theorem gives the entire sufficiently-small-viscosity
assertion in AIM Problem 415, with $\eps_T=1$. The time threshold obtained
below is sufficient; no optimality is asserted. The result concerns an
autonomous field with a strictly positive normal component. The case of a
merely nonnegative normal component, and the corresponding time-dependent
problem, are outside the statement.

The main choice is a solution of a Neumann Poisson problem with normal
derivative $b/2$. Conjugating by its exponential leaves a drift tangent to
$\Gamma$ and a positive Robin coefficient. The singular potential created by
the conjugation is included in the second part of the Carleman operator.
Its mixed product has the same admissible scale as the first-order terms.
The two weights agree on the boundary, so their tangential-gradient boundary
terms cancel, and the Robin substitution leaves a positive normal term. This
gives an observation bound of order $\exp(K/\eps)$ on a fixed interval.
Repeated flushing dissipation then dominates that exponential loss.

Section \ref{sec:adjoint} constructs the conjugation. Section
\ref{sec:observation} proves fixed-time observation, and Section
\ref{sec:uniform} combines it with dissipation and constructs the control.
The full mixed-product calculation and the justification for $L^2$ terminal
data are given in Appendices \ref{app:products} and \ref{app:regularity}.
All norms without a specified set are spatial $L^2(\Omega)$ norms. We work
over the reals; applying the argument to real and imaginary parts gives the
same bound for complex data.

\section{The adjoint and a boundary conjugation}\label{sec:adjoint}

Put $H=L^2(\Omega)$ and $V=H^1(\Omega)$, with $H$ identified as a subspace of
$V'$ by its inner product. Define
\[
 \mathfrak a(p,q)=\eps\int_\Omega\nabla p\cdot\nabla q
                        +\int_\Omega(B\cdot\nabla p)q.
\]
A weak solution of \eqref{state} belongs to
$L^2(0,T;V)\cap H^1(0,T;V')$ and satisfies
$\langle y_t,v\rangle+\mathfrak a(y,v)=(1_\omega h,v)$ for every $v\in V$.
The form and its transpose are bounded. Moreover,
\[
 \mathfrak a(v,v)\ge\frac\eps2\|\nabla v\|_2^2
                    -\frac{\|B\|_\infty^2}{2\eps}\|v\|_2^2.
\]
Adding a sufficiently large multiple $\kappa(v,v)$ makes both forms coercive.
Lions' variational existence theorem, in the form stated in
\cite[Theorem 1.1]{Dier}, gives existence, uniqueness and continuity in $H$.
The shift is removed by $v(t)=e^{-\kappa t}y(t)$, with the forcing multiplied
by $e^{-\kappa t}$. These are fixed-$\eps$ existence statements.

The backward adjoint with arbitrary terminal value $f\in H$ satisfies
$-\langle\phi_t,v\rangle+\mathfrak a(v,\phi)=0$. Its strong form is
\begin{equation}\label{adjoint}
 \begin{cases}
 \phi_t+\eps\Delta\phi+B\cdot\nabla\phi+(\divg B)\phi=0,\\
 \eps\partial_\nu\phi+b\phi=0,\\
 \phi(T)=f.
 \end{cases}
\end{equation}
Testing the variational equation and integrating the drift once gives
\begin{equation}\label{energy}
 \frac{d}{dt}\|\phi(t)\|_2^2
 =2\eps\|\nabla\phi(t)\|_2^2+\int_\Gamma b\phi(t)^2
                       -\int_\Omega(\divg B)\phi(t)^2
 \ge-M\|\phi(t)\|_2^2,
 \quad M=\|\divg B\|_\infty.
\end{equation}
Thus, for $t\le r$,
\begin{equation}\label{energy-transfer}
 \|\phi(t)\|_2^2\le e^{M(r-t)}\|\phi(r)\|_2^2.
\end{equation}
The time product rule in $L^2(V)\cap H^1(V')$ follows by time mollification;
both resulting products converge in $L^1$. This justifies the energy identity
for weak solutions.

To make the drift tangent on the boundary, set
\begin{equation}\label{poisson}
 k=\frac1{2|\Omega|}\int_\Gamma b,\qquad
 \Delta g=k,\qquad\partial_\nu g=b/2,\qquad\int_\Omega g=0.
\end{equation}
The data have the exact Neumann compatibility. Since $b\in C^{1,1}(\Gamma)$,
Nardi's theorem \cite[Theorem 3.1]{Nardi}, with $\alpha=1/2$, gives
$g\in C^{2,1/2}(\overline\Omega)$. Its boundary datum may be extended to the
closure through a smooth collar, as in that theorem's notation.

Define
\begin{equation}\label{gauge-coefficients}
 \begin{gathered}
 u=e^{g/\eps}\phi,\qquad D=B-2\nabla g,\qquad \beta=b/2,\\
 W=|\nabla g|^2-B\cdot\nabla g,\qquad c=\divg B-k.
 \end{gathered}
\end{equation}
Differentiating this multiplier in \eqref{adjoint} yields
\begin{equation}\label{gauged}
 u_t+\eps\Delta u+D\cdot\nabla u+(W/\eps+c)u=0,
 \qquad \eps\partial_\nu u+\beta u=0.
\end{equation}
We have $D\cdot\nu=0$, $D,W\in W^{1,\infty}(\Omega)$,
$c\in L^\infty(\Omega)$, and $\beta\ge c_0/2$.
Only the boundary normal derivative of $g$ was prescribed to cancel the normal
component of $D$. The choice of a constant Laplacian keeps $c$ bounded
without requiring third derivatives of $g$.

\section{Observation on a fixed interval}\label{sec:observation}

We first work on $Q=\Omega\times(0,1)$ and
$\Sigma=\Gamma\times(0,1)$. Choose $\omega'\Subset\omega$ and
$\eta\in C^2(\overline\Omega)$ such that
\begin{equation}\label{eta}
 \eta>0\ \text{in }\Omega,\quad \eta=0\ \text{on }\Gamma,
 \quad\|\eta\|_\infty=1,\quad
 |\nabla\eta|\ge\delta>0\ \text{on }\overline\Omega\setminus\omega'.
\end{equation}
The existence of this weight is the construction stated in
\cite[equation (64)]{Outward}. Write
$m=|\nabla\eta|^2$ and $a=\partial_\nu\eta$. On $\Gamma$,
$\nabla\eta=a\nu$, $m=a^2$, and $a\le-\delta$.

For $\sigma\in\{+1,-1\}$, define
\begin{equation}\label{weights}
 \xi_\sigma=\frac{e^{4\lambda+\sigma\lambda\eta}}{t(1-t)},\qquad
 \alpha_\sigma=\frac{e^{6\lambda}-e^{4\lambda+\sigma\lambda\eta}}{t(1-t)},
 \qquad w_\sigma=e^{-s\alpha_\sigma}u.
\end{equation}
Both weights agree on $\Sigma$, and so do $w_+$, $w_-$ and their time and
tangential derivatives. We denote their common boundary values by
$\xi,\alpha,w$. The opposite spatial derivatives of the two weights will
remove the tangential boundary norm.

\begin{proposition}\label{carleman}
There are a fixed $\lambda\ge1$ and constants $K,C>0$, depending only on the
fixed geometry and coefficients in \eqref{gauged}, such that every solution
of \eqref{gauged} on $Q$ satisfies
\begin{equation}\label{weighted-observation}
 \int_Qe^{-2s\alpha_+}\xi_+^3u^2\dd
 \le C\int_{\omega\times(0,1)}e^{-2s\alpha_+}\xi_+^3u^2\dd,
 \qquad s=K/\eps,\quad 0<\eps<1.
\end{equation}
The constants are independent of $\eps$.
\end{proposition}

\begin{proof}
We prove the estimate for the dense regular class described in Appendix
\ref{app:regularity}. All unlabelled integrals in this proof are over $Q$.
For a fixed $\sigma$, omit its subscript on $\xi,\alpha,w$, and put
$q=D\cdot\nabla\eta$. Split the transformed equation as
\begin{align}
 L_1w&=-2\eps s\lambda^2m\xi w
       -2\sigma\eps s\lambda\xi\nabla\eta\cdot\nabla w
       +w_t+D\cdot\nabla w,\label{split1}\\
 L_2w&=\eps s^2\lambda^2m\xi^2w+\eps\Delta w
       +s\alpha_tw-\sigma s\lambda\xi q w+(W/\eps)w,\label{split2}\\
 L_3w&=\sigma\eps s\lambda\Delta\eta\,\xi w
       -\eps s\lambda^2m\xi w-cw.\label{split3}
\end{align}
Indeed, $\nabla\alpha=-\sigma\lambda\xi\nabla\eta$ and
$\Delta\alpha=-\lambda^2m\xi-\sigma\lambda\Delta\eta\,\xi$ imply
$L_1w+L_2w=L_3w$. Hence
\begin{equation}\label{square-identity}
 \|L_1w\|_{L^2(Q)}^2+\|L_2w\|_{L^2(Q)}^2+2(L_1w,L_2w)_{L^2(Q)}
                  =\|L_3w\|_{L^2(Q)}^2.
\end{equation}
The two products generating the leading mass contribute coefficients
$-2$ and $+3$, leaving
\begin{equation}\label{leading}
 \eps^2s^3\lambda^4\int m^2\xi^3w^2
       +\eps^2s\lambda^2\int m\xi|\nabla w|^2,
\end{equation}
together with a nonnegative directional-gradient square. Appendix
\ref{app:products} gives every mixed product. Before the boundary condition
is imposed, their boundary contribution, apart from the potential, is
\begin{align}
 I_\sigma=\int_\Sigma\bigl[&-\sigma\eps^2s^3\lambda^3ma\xi^3w^2
 -2\eps^2s\lambda^2m\xi w d_\sigma
 -\sigma\eps^2s\lambda a\xi d_\sigma^2\notag\\
 &+\sigma\eps^2s\lambda a\xi|\nabla_\Gamma w|^2
 +\eps(w_t+D\cdot\nabla_\Gamma w)d_\sigma
 -\sigma\eps s^2\lambda\alpha_t\xi a w^2\bigr]\ds,
 \label{boundary-unsummed}
\end{align}
where $d_\sigma=\partial_\nu w_\sigma$.

The singular potential was retained in $L_2$. Since $W$ is autonomous and
$D\cdot\nu=0$, its mixed product is exactly
\begin{align}
 (L_1w,(W/\eps)w)=
 \int\bigl[&-s\lambda^2Wm\xi
 +\sigma s\lambda\xi(\nabla W\cdot\nabla\eta+W\Delta\eta)
 -(2\eps)^{-1}\divg(WD)\bigr]w^2\dd\notag\\
 &-\sigma s\lambda\int_\Sigma W\xi a w^2\ds.
 \label{potential-product}
\end{align}
In particular, its bulk terms have absolute value at most
$C\int(s\lambda^2\xi+1/\eps)w^2$. They do not create a squared potential
on the right of \eqref{square-identity}. The boundary terms in
\eqref{potential-product} cancel between the two signs.

The Robin condition in \eqref{gauged} gives
\begin{equation}\label{normal-weighted}
 d_\sigma=(-\beta/\eps+\sigma s\lambda\xi a)w.
\end{equation}
Using $m=a^2$ on $\Gamma$, substitute this into
\eqref{boundary-unsummed} and sum. Terms odd in $\sigma$ cancel. The scalar
normal term gives $4\eps s\lambda^2\beta a^2\xi w^2$; the difference of
the normal squares gives $4\eps s^2\lambda^2\beta a^2\xi^2w^2$.
The time term is $-2\beta w_tw$, whose integral vanishes.
The remaining drift term is $-\beta D\cdot\nabla_\Gamma(w^2)$, and
integration on the closed boundary gives
\begin{equation}\label{boundary-positive}
 I_++I_-=
 \int_\Sigma\left[4\eps s\lambda^2\beta a^2\xi
       +4\eps s^2\lambda^2\beta a^2\xi^2
       +\divg_\Gamma(\beta D)\right]w^2\ds.
\end{equation}
The divergence in this formula is bounded. Since $\beta\ge c_0/2$,
$a^2\ge\delta^2$ and $\xi$ has a positive lower bound, increasing $K$
makes \eqref{boundary-positive} nonnegative whenever $s\ge K/\eps$.
There is no remaining tangential-gradient boundary integral.

We give the bulk absorption explicitly. Apart from
\eqref{potential-product}, the absolute values of the errors are bounded by
\begin{align}
 C\int\bigl[&\eps^2s^3\lambda^3\xi^3+
 \eps^2s^2\lambda^4\xi^2+
 \eps s^2\lambda^3\xi^2+
 \eps s^2\lambda^2\xi^3\notag\\
 &+s\xi^3+s\lambda\xi^2+s\lambda^2\xi\bigr]w^2
 +C\int[\eps^2s\lambda\xi+\eps^2\lambda^2+\eps]|\nabla w|^2.
 \label{errors}
\end{align}
The terms with $s^3\lambda^3$ and $s\lambda$ have constants independent of
$\lambda$ and are absorbed outside $\omega'$ by first choosing $\lambda$
large. Now fix $\lambda$. The estimates
\[
 |\xi_t|\le C\xi^2,\quad |\alpha_t|\le C\xi^2,
 \quad |\alpha_{tt}|\le C\xi^3,\quad
 |\nabla\alpha_t|\le C\lambda\xi^2
\]
and the positive lower bound on $\xi$ show that all the other ratios of
\eqref{errors} to \eqref{leading} tend to zero as $K\to\infty$ when
$s\ge K/\eps$. The terms $\eps s^2\xi^3$ and $s\xi^3$, for example, have
ratios at most $C/K$ and $C/K^2$. The ratios from
\eqref{potential-product} are at most $C/K^2+C/K^3$. Also
\[
 \|L_3w\|_{L^2(Q)}^2
 \le C\int(\eps^2s^2\lambda^4\xi^2+\|c\|_\infty^2)w^2,
\]
which is absorbed at the same scale outside $\omega'$. The first-order
cross term involving $\nabla(m\xi)$ is estimated after just one integration
by parts, so only two derivatives of $\eta$ are needed.

Set, with full-cylinder integrals,
\begin{align*}
 P_\sigma&=\eps^2s^3\lambda^4\int\xi_\sigma^3w_\sigma^2
           +\eps^2s\lambda^2\int\xi_\sigma|\nabla w_\sigma|^2,\\
 N_\sigma&=\|L_1w_\sigma\|_{L^2(Q)}^2+\|L_2w_\sigma\|_{L^2(Q)}^2.
\end{align*}
Since $m\ge\delta^2$ outside $\omega'$, the preceding calculations imply
\begin{align}\label{local-gradient}
 \sum_\sigma(N_\sigma+P_\sigma)
 \le C\sum_\sigma\biggl[&\eps^2s^3\lambda^4
        \int_{\omega'\times(0,1)}\xi_\sigma^3w_\sigma^2
 +\eps^2s\lambda^2\int_{\omega'\times(0,1)}
                         \xi_\sigma|\nabla w_\sigma|^2\biggr].
\end{align}
The added mass and gradient inside $\omega'$ are included on the right.

To remove the local gradient, recover the Laplacian from $L_2$:
\begin{equation}\label{laplacian}
 H_\sigma:=\eps^2\int(s\xi_\sigma)^{-1}|\Delta w_\sigma|^2
                \le C(N_\sigma+P_\sigma).
\end{equation}
Indeed, after squaring \eqref{split2} and dividing by $s\xi_\sigma$,
the term with $s^2\xi_\sigma^2$ is bounded by the mass in $P_\sigma$;
the $\alpha_t$ and $q$ terms have ratio at most $C/K^2$ to that mass.
The potential term has ratio
$C/(\eps^4s^4\xi_\sigma^4)\le C/K^4$.
Take $\theta\in C_c^\infty(\omega)$ equal to one near
$\overline{\omega'}$. The interior identity
\[
 \int\theta^2\xi_\sigma|\nabla w_\sigma|^2
 =-\int\theta^2\xi_\sigma w_\sigma\Delta w_\sigma
   -\int\nabla(\theta^2\xi_\sigma)\cdot\nabla w_\sigma\,w_\sigma
\]
and Young's inequality give, for every $\rho>0$,
\begin{equation}\label{cutoff}
 \eps^2s\lambda^2\int_{\omega'\times(0,1)}
                       \xi_\sigma|\nabla w_\sigma|^2
 \le\rho H_\sigma+C_\rho\eps^2s^3\lambda^4
                  \int_{\omega\times(0,1)}\xi_\sigma^3w_\sigma^2.
\end{equation}
Here one first absorbs half the cutoff-gradient integral and then pairs
$\Delta w_\sigma$ with $w_\sigma$ using weights
$(s\xi_\sigma)^{-1}$ and $(s\xi_\sigma)^3$.
Choose $\rho$ small in \eqref{cutoff}, use \eqref{laplacian}, and absorb
into \eqref{local-gradient}. Retain the positive mass and use
$\xi_-\le\xi_+$, $\alpha_-\ge\alpha_+$ to get
\eqref{weighted-observation}. Appendix \ref{app:regularity} passes this
estimate to every $L^2$ terminal datum.
\end{proof}

\begin{proposition}\label{coarse}
There are constants $C_o,K_o>0$, independent of $\eps\in(0,1)$, such that
every adjoint \eqref{adjoint} on a translated unit interval satisfies
\begin{equation}\label{coarse-observation}
 \|\phi(t)\|_2^2\le C_oe^{K_o/\eps}
              \int_t^{t+1}\|\phi(r)\|_{L^2(\omega)}^2\,dr.
\end{equation}
\end{proposition}
\begin{proof}
On $[1/3,2/3]\times\overline\Omega$ the weight $\alpha_+$ is bounded
above and $\xi_+$ bounded below. On all of $Q$, the weight
$e^{-2s\alpha_+}\xi_+^3$ is bounded above independently of $s\ge K\ge1$:
the exponential bounds every power of $[t(1-t)]^{-1}$. Thus
Proposition \ref{carleman} gives
\[
 \int_{1/3}^{2/3}\|u(t)\|_2^2\,dt
 \le C e^{C/\eps}\int_0^1\|u(t)\|_{L^2(\omega)}^2\,dt.
\]
Undoing $u=e^{g/\eps}\phi$ adds at most $2\osc(g)/\eps$ to the exponent.
Equation \eqref{energy-transfer} transfers middle-time mass to the initial
time. Finally, autonomy permits translation of the unit interval.
\end{proof}

\section{Uniform observation and the control}\label{sec:uniform}

We return to the original adjoint \eqref{adjoint}. The general-drift
dissipation estimate \cite[Proposition 4.7(1), equation (51)]{Outward} has the
following form. If $B\in L^\infty_tW^{1,\infty}_x$ has nonnegative normal
component and satisfies the source's flushing condition on an ambient
interval $[0,L]$, with $T_0<L$, then
\begin{equation}\label{dissipation}
 \|\phi(t_0-T_0)\|_2^2\le C_d\left[
 e^{-\delta_d/\eps}\|\phi(t_0)\|_2^2+
 \int_{t_0-T_0}^{t_0}\|\phi(t)\|_{L^2(\omega)}^2\,dt\right],
 \quad t_0\in[T_0,L].
\end{equation}
This applies to weak adjoints with arbitrary $L^2$ terminal data.
$C_d$ is independent of $\eps$ and $t_0$. The decay exponent
$\delta_d>0$, specified in \cite[Proposition 4.5]{Outward}, depends on the
flushing radius, $T_0$ and the $L^\infty_tW^{1,\infty}_x$ drift norm.
For the autonomous field this exponent can be fixed before choosing $L$.
The prefactor $C_d$ is allowed to depend on the subsequently fixed $L$.

To check the source condition, its flow $\Phi(t,t_0,x)$ solves the forward
equation for $B$ with value $x$ at $t_0$. Autonomy gives
\[
 \Phi(t,t_0,x)=F_{t_0-t}(x).
\]
For any $x\in\overline\Omega$ and $t_0\ge T_0$, choose
$t=t_0-s_x\in(t_0-T_0,t_0)$. Equation \eqref{flushing} implies
\[
 \Phi\bigl(t,t_0,\overline{\mathbb B}(x,r_0/2)\bigr)\subset\omega.
\]
This is its closed-ball flushing condition. The source's boundary assumption
follows from \eqref{outward}. No gradient assumption occurs in
\eqref{dissipation}.

\begin{proof}[Proof of Theorem \ref{main}]
Choose an integer $m\ge1$ such that $m\delta_d>K_o$, with $K_o$ from
\eqref{coarse-observation}, and set $L=mT_0+1$. Fix the dissipation
prefactor on this ambient interval and increase it to $C_d\ge1$.
Iterating \eqref{dissipation} on the $m$ successive blocks gives
\begin{equation}\label{iterated}
 \|\phi(0)\|_2^2\le C_d^m\left[
 e^{-m\delta_d/\eps}\|\phi(mT_0)\|_2^2+
                     \int_0^{mT_0}\|\phi(t)\|_{L^2(\omega)}^2\,dt\right].
\end{equation}
Apply \eqref{coarse-observation} on $[mT_0,L]$. Since
$m\delta_d-K_o>0$, \eqref{iterated} yields
\[
 \|\phi(0)\|_2^2\le C_*\int_0^L\|\phi(t)\|_{L^2(\omega)}^2\,dt,
 \qquad C_*=C_d^m\max\{1,C_o\}.
\]
For $T\ge L$, restrict its adjoint to $[0,L]$ with terminal datum
$\phi(L)$. Uniqueness identifies this restriction with the weak adjoint
just used. Consequently
\begin{equation}\label{uniform-observation}
 \|\phi(0)\|_2^2\le C_*\int_0^T\|\phi(t)\|_{L^2(\omega)}^2\,dt,
 \qquad T\ge L,\quad 0<\eps<1.
\end{equation}

Let $\mathcal M$ be the closure in $L^2((0,T)\times\omega)$ of these
adjoint observations. Equation \eqref{uniform-observation} makes
\[
 \ell(\phi|_\omega)=-(y_0,\phi(0))
\]
a well-defined bounded functional, of norm at most
$C_*^{1/2}\|y_0\|_2$. Extend it by continuity to $\mathcal M$.
Its Hilbert representative $h\in\mathcal M$ has the same norm bound.
One can construct it directly: if $\ell\ne0$, minimize the norm on
$\{z\in\mathcal M:\ell(z)=1\}$. The parallelogram identity makes a
minimizing sequence Cauchy. Its limit $z$ is orthogonal to $\ker\ell$,
so $h=z/\|z\|^2$ represents $\ell$. If $\ell=0$, take $h=0$.

For the state with this control, the weak time-product rule and the two
transposed forms give
\[
 (y(T),f)-(y_0,\phi(0))=\int_0^T\int_\omega h\phi.
\]
The right side equals $-(y_0,\phi(0))$. Therefore $(y(T),f)=0$ for every
$f\in L^2(\Omega)$, and $y(T)=0$. Taking $T_*=L$ and $C=C_*^{1/2}$ proves
all the asserted quantifiers.
\end{proof}

\appendix
\section{The mixed products}\label{app:products}

Here all bulk integrals are over $Q$. Fix $\sigma$ and use its $\xi,\alpha,w$;
write $v=\nabla\eta$, $m=|v|^2$, $q=D\cdot v$ and $d=\partial_\nu w$.
To group the calculation, denote the four terms of $L_1$ in
\eqref{split1} by
\[
 A=-2\eps s\lambda^2m\xi w,\quad
 G=-2\sigma\eps s\lambda\xi v\cdot\nabla w,\quad
 J=w_t,\quad Z=D\cdot\nabla w.
\]
Denote the first four terms of $L_2$ by
\[
 R=\eps s^2\lambda^2m\xi^2w,\quad H_0=\eps\Delta w,
 \quad S=s\alpha_tw,\quad U=-\sigma s\lambda\xi q w.
\]
These letters are local to this appendix. Integration by parts gives
\begin{align}
 (A,R)&=-2\eps^2s^3\lambda^4\int m^2\xi^3w^2,\notag\\
 (G,R)&=3\eps^2s^3\lambda^4\int m^2\xi^3w^2
       +\sigma\eps^2s^3\lambda^3\int\divg(mv)\xi^3w^2\notag\\
       &\hspace{1cm}-\sigma\eps^2s^3\lambda^3\int_\Sigma ma\xi^3w^2,\notag\\
 (J,R)&=-\eps s^2\lambda^2\int m\xi\xi_t w^2,\notag\\
 (Z,R)&=-\frac{\eps s^2\lambda^2}{2}\int\divg(m\xi^2D)w^2.
 \label{products-R}
\end{align}
In the second line one uses $\nabla\xi=\sigma\lambda\xi v$, and the drift
has no boundary term because $D\cdot\nu=0$.

The products against the Laplacian are
\begin{align}
 (A,H_0)&=2\eps^2s\lambda^2\int m\xi|\nabla w|^2
       +2\eps^2s\lambda^2\int\nabla(m\xi)\cdot\nabla w\,w\notag\\
       &\hspace{1cm}-2\eps^2s\lambda^2\int_\Sigma m\xi wd,\notag\\
 (G,H_0)&=2\sigma\eps^2s\lambda\int\xi D^2\eta(\nabla w,\nabla w)
       -\sigma\eps^2s\lambda\int\xi\Delta\eta|\nabla w|^2\notag\\
       &\quad-\eps^2s\lambda^2\int m\xi|\nabla w|^2
       +2\eps^2s\lambda^2\int\xi(v\cdot\nabla w)^2\notag\\
       &\quad-\sigma\eps^2s\lambda\int_\Sigma a\xi d^2
       +\sigma\eps^2s\lambda\int_\Sigma a\xi|\nabla_\Gamma w|^2,\notag\\
 (J,H_0)&=\eps\int_\Sigma w_t d,\notag\\
 (Z,H_0)&=\eps\int_\Sigma(D\cdot\nabla_\Gamma w)d
       +\frac\eps2\int(\divg D)|\nabla w|^2\notag\\
       &\hspace{1cm}-\eps\int\sum_{i,j}(\partial_iD_j)(\partial_iw)(\partial_jw).
 \label{products-H}
\end{align}
The cross term in the first line is bounded by
\[
 C\eps^2\lambda^2\int|\nabla w|^2+
 C\eps^2s^2\lambda^4\int\xi^2w^2.
\]
For the second line, apply the divergence identity to
$\xi v\cdot\nabla(|\nabla w|^2)$ after the first integration by parts.
This avoids differentiating $\eta$ three times.

The time-weight products are
\begin{align}
 (A,S)&=-2\eps s^2\lambda^2\int m\xi\alpha_tw^2,\notag\\
 (G,S)&=\sigma\eps s^2\lambda\int\divg(\xi\alpha_tv)w^2
       -\sigma\eps s^2\lambda\int_\Sigma\xi\alpha_ta w^2,\notag\\
 (J,S)&=-\frac s2\int\alpha_{tt}w^2,\notag\\
 (Z,S)&=-\frac s2\int\divg(\alpha_tD)w^2.
 \label{products-S}
\end{align}
The remaining first-order products are
\begin{align}
 (A,U)&=2\sigma\eps s^2\lambda^3\int m\xi^2q w^2,\notag\\
 (G,U)&=-\eps s^2\lambda^2\int\divg(\xi^2qv)w^2,\notag\\
 (J,U)&=\frac{\sigma s\lambda}{2}\int\xi_tq w^2,\notag\\
 (Z,U)&=\frac{\sigma s\lambda}{2}\int\divg(\xi qD)w^2.
 \label{products-U}
\end{align}
Their boundary terms vanish because $q=0$ and $D\cdot\nu=0$ on $\Gamma$.
The last four products, against $(W/\eps)w$, sum to
\eqref{potential-product}. The time product is zero, while the other three
are the integrations of the scalar term, of
$-\sigma s\lambda W\xi v\cdot\nabla(w^2)$, and of
$(2\eps)^{-1}WD\cdot\nabla(w^2)$. These are all twenty mixed products.
The integration order is the tangential-drift calculation of
\cite[Section 6]{Tangential}, with the potential retained and the Robin
substitution performed in \eqref{normal-weighted}--\eqref{boundary-positive}.
No Neumann cancellation from that source is assumed here.

\section{Justification for rough terminal data}\label{app:regularity}

We give the regularity and density argument at a fixed $\eps>0$. On a smooth
bounded domain, normal collar coordinates and the fundamental theorem of
calculus give, for $z\in H^1(\Omega)$ and every $\theta>0$,
\begin{equation}\label{trace}
 \|z\|_{L^2(\Gamma)}^2\le\theta\|\nabla z\|_2^2+C_\theta\|z\|_2^2.
\end{equation}
For smooth $z$, average its boundary squared value over a collar and estimate
the integral of $2|z\partial_rz|$ by Cauchy--Schwarz and Young. The collar
Jacobian is bounded above and below. Approximation in boundary graph charts
passes the estimate to $H^1$.

For a smooth Neumann function $z$, integrations by parts yield
\[
 \int_\Omega(\Delta z)^2-\int_\Omega|D^2z|^2
       =\int_\Gamma\mathrm{II}(\nabla_\Gamma z,\nabla_\Gamma z),
 \qquad\mathrm{II}(X,Y)=(\nabla_X\nu)\cdot Y.
\]
Before using $\partial_\nu z=0$, the boundary expression is
$\Delta z\partial_\nu z-D^2z(\nabla z,\nu)$.
Differentiating the normal condition tangentially gives the displayed
curvature term. Bound its absolute value by
$C\|\nabla z\|_{L^2(\Gamma)}^2$ and use \eqref{trace} on each derivative
with $\theta$ small. This gives
\begin{equation}\label{H2}
 \|z\|_{H^2}\le C(\|\Delta z\|_2+\|z\|_{H^1}).
\end{equation}
It also holds for $C^{2,1/2}$ Neumann functions. Even extension in the normal
coordinate and even mollification approximate them in $C^2$ while preserving
the normal condition; a finite partition with zero normal derivative
combines the local approximations.

If $z\in H^1$ is a weak Neumann function with $\Delta z\in L^2$, put
$f=-\Delta z+z$ and approximate $f$ by smooth $f_j$ in $L^2$.
Nardi's theorem \cite[Theorem 5.2]{Nardi} provides $C^{2,1/2}$ solutions
$z_j$ of $-\Delta z_j+z_j=f_j$, $\partial_\nu z_j=0$.
Their variational identity gives
$\|z_j-z_l\|_{H^1}\le\|f_j-f_l\|_2$; their Laplacians are Cauchy in $L^2$,
so \eqref{H2} makes them Cauchy in $H^2$. The limit equals $z$, since
testing the homogeneous $-\Delta+1$ equation proves uniqueness.
Thus \eqref{H2} holds for weak Neumann functions.

For a weak Robin function $v$ with $\Delta v\in L^2$ and
$\eps\partial_\nu v+bv=0$, the multiplier $e^{2g/\eps}v$ is Neumann and
has an $L^2$ Laplacian. Its multiplier and inverse preserve $H^2$. Hence
\begin{equation}\label{Robin-H2}
 \|v\|_{H^2}\le C_\eps(\|\Delta v\|_2+\|v\|_{H^1}).
\end{equation}
This constant is used only at fixed $\eps$.

Let $S(r)f$ denote the original adjoint after reversing time. Weak uniqueness
gives its semigroup property; \eqref{energy-transfer} gives
$\|S(r)\|\le e^{Mr/2}$, and the continuous initial trace gives strong
continuity. Choose nonnegative
$\rho_n\in C_c^\infty((0,1/n))$, of integral one, and set
$f_n=\int\rho_n(r)S(r)f\,dr$. Then $f_n\to f$ in $L^2$ and
\begin{equation}\label{time-smoothing}
 \partial_q^j S(q)f_n=(-1)^j\int\rho_n^{(j)}(r)S(q+r)f\,dr.
\end{equation}
This formula differentiates the translated compact test function, rather
than presupposing differentiability of the orbit. Its initial values belong
to $H^1$ because weak solutions are $L^2$ in time with values in $H^1$.
The fixed-$\eps$ strong-solution and continuity results
\cite[Propositions 3.6 and 3.9]{Outward} make each of these orbits continuous
in $H^1$. The equation then makes its Laplacian continuous in $L^2$;
\eqref{Robin-H2} makes it, and each of its time derivatives, continuous in
$H^2$. Conjugation by $e^{g/\eps}$ preserves this regularity. The Sobolev
Green identities and their boundary traces therefore justify every
integration in Section \ref{sec:observation} and Appendix \ref{app:products}.
The time endpoint weights vanish faster than every polynomial, so there
are no temporal endpoint terms.

Finally \eqref{energy-transfer} gives convergence of the approximating
adjoints in $C([0,T];L^2)$ and hence in the observation norm. At each fixed
$\eps$ the gauge is a bounded invertible multiplier. The weighted estimates
and \eqref{coarse-observation} pass to the limit with their already obtained
constants. This passage is made at fixed $\eps$ and does not interchange a
data-density limit with the vanishing-viscosity limit.

\section*{AI-use statement}
OpenAI Codex provided substantial assistance in developing the argument,
checking its calculations, locating and comparing the cited sources, writing
the manuscript, and preparing the submission files. The argument and its
conclusions are submitted for mathematical review.

\begin{thebibliography}{9}
\bibitem{Barcena}
J.~A. B\'arcena-Petisco,
\emph{Cost of null controllability for parabolic equations with vanishing
diffusivity and a transport term}, ESAIM: Control, Optimisation and Calculus
of Variations \textbf{27} (2021), article 106.
\href{https://doi.org/10.1051/cocv/2021103}{doi:10.1051/cocv/2021103}.
Source text: \href{https://hal.science/hal-02455632v4}{HAL hal-02455632v4}.

\bibitem{Dier}
D. Dier, \emph{Non-Autonomous Maximal Regularity for Forms of Bounded
Variation}, preprint, \href{https://arxiv.org/abs/1406.2884v2}{arXiv:1406.2884v2}
(2014).

\bibitem{Tangential}
F. Et-Tahri, J.~A. B\'arcena-Petisco, I. Boutaayamou and L. Maniar,
\emph{On uniform null-controllability of tangential transport-diffusion
equations with vanishing viscosity limit}, preprint,
\href{https://arxiv.org/abs/2308.07497v1}{arXiv:2308.07497v1} (2023).

\bibitem{Outward}
F. Et-Tahri, J.~A. B\'arcena-Petisco, I. Boutaayamou and L. Maniar,
\emph{On uniform null controllability of transport-diffusion equations with
vanishing viscosity limit}, Mathematical Methods in the Applied Sciences
\textbf{48} (2025), 6531--6552.
\href{https://doi.org/10.1002/mma.10693}{doi:10.1002/mma.10693}.
The theorem and equation locators in this paper refer to
\href{https://arxiv.org/abs/2412.16660v1}{arXiv:2412.16660v1} (2024).

\bibitem{Nardi}
G. Nardi, \emph{Schauder estimate for solutions of Poisson's equation with
Neumann boundary condition}, L'Enseignement Math\'ematique \textbf{60} (2014),
421--435. \href{https://doi.org/10.4171/LEM/60-3/4-9}{doi:10.4171/LEM/60-3/4-9}.
The theorem locators refer to
\href{https://arxiv.org/abs/1302.4103v4}{arXiv:1302.4103v4} (2015).
\end{thebibliography}
\end{document}
